\chapter{Conclusão}
\label{cap:conclusao}

\section{Síntese das Contribuições}

O trabalho propôs e implementou um Data Warehouse para análise da execução orçamentária de Juiz de Fora, com o duplo objetivo de tornar dados públicos analiticamente acessíveis e de oferecer um modelo replicável para outros municípios. A revisão dos objetivos específicos, contra o que foi efetivamente entregue, ajuda a fixar o que esse trabalho contribui.

O primeiro objetivo previa a coleta e documentação de dados de receitas e despesas do Portal da Transparência da Prefeitura. Foi atendido: o \textit{pipeline} de extração consome quatro fontes (despesa mensal consolidada, receita mensal comparativa, receita mensal prevista e ementário de natureza de receita da STN), com tratamento explícito de heterogeneidades como mudança de \textit{layout} no meio do histórico e cabeçalho de planilha sem posição fixa. As inconsistências encontradas foram tratadas como achados de pesquisa e documentadas no Capítulo~\ref{cap:dw}.

O segundo objetivo previa a modelagem dimensional segundo Kimball. Foi atendido. O \textit{warehouse} conta com oito dimensões e três tabelas de fatos organizadas em esquema estrela compartilhando \texttt{dim\_tempo}. Cada decisão de granularidade foi escolhida em função do que a fonte de fato publica, sem inventar atomicidade ausente nos dados originais.

O terceiro objetivo previa a implementação de processos de ETL. Foi atendido com Dagster como orquestrador, Python e \texttt{pandas} para extração, DuckDB como banco analítico e dbt para modelagem em SQL. A unidade de idempotência foi definida como partição temporal e três mecanismos garantem que a reexecução da mesma partição produza o mesmo estado final do \textit{warehouse}.

O quarto objetivo previa a construção de \textit{dashboards} interativos. Foi atendido. Sete páginas analíticas em Streamlit consomem o \textit{warehouse} e respondem perguntas como execução orçamentária por área, evolução temporal de receitas, concentração de fornecedores e desvio entre previsão e arrecadação. \comment[inline]{Quando o Capítulo~\ref{cap:resultados} estiver escrito, retomar aqui os indicadores mais relevantes observados nos dados de Juiz de Fora.}

O quinto objetivo previa a avaliação de usabilidade. \comment[inline]{Sintetizar os achados do Capítulo~\ref{cap:avaliacao} aqui, com base na comparação entre as visualizações produzidas e o formato original do portal.}

Além dos objetivos declarados, o trabalho contribui de três formas adicionais. Primeiro, oferece uma referência arquitetural concreta para projetos de DW municipal, com escolhas técnicas explicitadas e justificadas. Segundo, documenta de forma estruturada as limitações da fonte (ausência de orçamento autorizado, granularidade só mensal, classificação MCASP incompleta), que costumam ser tratadas como ruído silencioso em projetos similares. Terceiro, articula a literatura técnica de Kimball com a tradição freireana de educação libertadora, indicando que a arquitetura de dados não é decisão neutra.

\section{Considerações Finais}

O processo de construção deste \textit{warehouse} ensinou algo que não estava previsto no plano inicial. O problema central da transparência orçamentária municipal não é técnico no sentido mais estreito da palavra. As tecnologias necessárias existem e estão maduras. DuckDB, dbt, Dagster, Streamlit, Python: tudo isso é gratuito, bem documentado e relativamente fácil de operar. O problema é de outra ordem.

O problema está na distância entre o que o portal publica e o que o cidadão precisa para fiscalizar. O portal publica relatórios mensais consolidados em planilhas Excel com cabeçalhos sem padronização. O cidadão precisa de séries temporais comparáveis, classificações inteligíveis e indicadores que dialoguem com sua experiência cotidiana. Reduzir essa distância é trabalho de arquitetura de dados, e arquitetura de dados é, em última leitura, uma decisão sobre quem deve conseguir interpretar a informação pública. Quando se decide manter o dado em formato consolidado mensal, sem detalhe transacional, decide-se que o cidadão verá o agregado e não o ato. Quando se publica em Excel com cabeçalho de duas linhas e códigos sem ementário, decide-se que apenas quem sabe \textit{parsing} chega aos números. Essas decisões podem ser revertidas. É escolha política, não restrição tecnológica.

A leitura freireana ganha contorno técnico nesse ponto. Conscientização orçamentária, traduzida para o trabalho de quem constrói \textit{data warehouses}, significa modelar o dado pensando em quem vai lê-lo, e não apenas em quem vai gerá-lo. O Kimballiano ``orientado a assunto'' encontra, nesse encontro, sua versão democrática.

\section{Trabalhos Futuros}

A continuação natural do trabalho segue cinco direções concretas, ordenadas por relevância para o controle social:

(1) Integração do orçamento autorizado pela LOA, viabilizando indicadores de percentual de execução sobre o autorizado, que são a métrica padrão usada por órgãos de controle.

(2) Decomposição da funcional programática em dimensões separadas para função, subfunção, programa e ação, com base no ementário oficial da STN. Essa decomposição substitui o dicionário ad hoc de classificação de funções hoje presente no código do \textit{dashboard}.

(3) Captura de licitações e contratos como processos de negócio adicionais, permitindo rastrear a cadeia completa do gasto (de licitação a pagamento) e analisar tempo médio entre etapas.

(4) Implementação de SCDs tipo 2 nas dimensões sujeitas a mudança histórica (em particular reorganizações administrativas), permitindo análises ``como era'' \textit{versus} ``como é''.

(5) Replicação do modelo em pelo menos um município de porte semelhante, com análise comparativa, para validar a generalidade do esquema dimensional proposto e identificar pontos de adaptação local.

Cada uma dessas direções é viável com o ferramental atual e não exige reescrita das camadas já implementadas. Trabalho futuro mais ambicioso inclui automação de coleta com agendamento real (hoje as partições suportam \textit{backfill} sob demanda, mas não há \textit{scheduler} ativo), exposição do \textit{warehouse} como serviço público acessível e avaliação empírica de impacto em ações concretas de controle social.

O \textit{warehouse} construído aqui, tomado em conjunto com a discussão teórica que o sustenta, não esgota o problema da transparência orçamentária. Oferece, em vez disso, uma peça de evidência de que é possível, em escopo de um TCC, atravessar o caminho entre o dado bruto que a prefeitura publica e a pergunta que o cidadão de fato faria. Esse atravessamento é a contribuição.
